% !TEX root = ~/Code/TN_LOWE_RECONSTRUCTION/main.tex
\section{Introduction} \label{sec:INTRODUCTION}
%There is a long history of advancements in particle and astroparticle physics resulting from measuring astrophysical neutrinos with energies in the \SI{10}{\mega\electronvolt} range~\cite{Ahmad2001,Abe2016}. The sun produces a continuous flux of roughly \SI{7e10}{\per\second\per\centi\meter\squared} neutrinos as a byproduct of fusion reactions within the stellar core~\cite{Vinyoles_2018}. Despite the confirmation of solar neutrino oscillations as a solution to the \textit{Solar Neutrino Problem}~\cite{PhysRevLett.65.2233}, several questions remain~\cite{Gann2021ndb}. 
The Deep Underground Neutrino Experiment, as a next--generation large--scale neutrino detector, will have the ability to shed light on some missing pieces of the solar neutrino puzzle. In particular, \ac{DUNE} will be sensitive to the two most energetic components of the solar neutrino flux, $^{8}\text{B}$ and $hep$ neutrinos, as well as, neutrinos from a nearby occurring supernovae. The high cross--section of \ac{CC} interaction on argon nuclei~\cite{abi2021supernova}, together with \ac{DUNE}'s large active mass and excellent energy resolution, will allow for precision measurements of the astrophysical neutrino flux and their energy spectra. These measurements will provide insights into neutrino properties (such as the MSW effect and potential non-standard interactions), solar physics (including metallicity and core temperature), and potentially new physics beyond the Standard Model.
% Solar neutrinos have been intensely studied, beginning over 50 years ago with the Homestake experiment, which observed only 39$\%$ of the neutrinos predicted by the standard solar model (SSM). Neutrino oscillations were soon postulated as a solution but were not observed until the SNO experiment measured the $\nu_e$ and total neutrino fluxes simultaneously, demonstrating flavour transitions.

\subsection{DUNE Far Detectors}
\ac{DUNE} will be made of four \ac{FD} modules based on the \ac{LArTPC} technology in the \ac{SURF} to survey a wide portfolio of neutrino physics topics\iffalse. The site has \SI{4700}{\metersofwater} of overburden which significantly reduces the rate of cosmic backgrounds throughout the detector and allows searches for low-energy astroparticle measurements.  The PIP-II accelerator upgrade at the \ac{FERMILab} will give a new neutrino beam directed at SURF to study long-baseline neutrino oscillations with \ac{DUNE}\fi. The \ac{FD}s will be built in two phases. The first stage will deploy two modules with approximately \SI{14}{\kilo\tonne} of active mass. Two distinct designs (or geometries), are going to be implemented: the \ac{HD} and the \ac{VD} designs (see \cref{fig:DUNE_FAR_DETECTORS}), named after the orientation of their drift fields. The \ac{HD}'s fiducial volume will extend \SI{14.4}{\meter} in the drift direction, aligned along the drift coordinate ($x$), \SI{12}{\meter} along the vertical coordinate $y$, and \SI{58}{\meter} along the beam--aligned coordinate ($z$). This module will be separated into four distinct drift volumes alternating three \ac{APA} and two \ac{CPA} sets, which are set to a drift potential of \SI{180}{\kilo\volt}, producing an electric field of \SI{500}{\volt\per\centi\meter}.
%Each APA is outfitted with three charge readout planes (CRP), consisting of a wire mesh, to monitor the charge deposited within the active argon volume.

\begin{figure}[ht]
    \centering
    \includegraphics[width=.49\linewidth]{Images/fd10kt.png}
    \includegraphics[width=.49\linewidth]{Images/fdvd10kt.png}
  \caption[width=1\linewidth]{(left) \ac{DUNE} \ac{FD} \ac{HD} module. (right) \ac{DUNE} \ac{FD} \ac{VD} module.\label{fig:DUNE_FAR_DETECTORS}}
\end{figure}

The \ac{FD} \ac{VD} module simplifies the traditional \ac{HD} design by dividing the active mass into just two drift volumes (of dimensions \numproduct{13.5 x 6.5 x 60}\unit{\metre\cubed}) which separated horizontally by the cathode (denoted simply as Cathode Plane). The equivalent of the \ac{CPA} in the \ac{VD} volume is called the \ac{CRP} and, differently to the thin and free frame--tensioned wires of a traditional \ac{LArTPC}, its building block is a perforated \ac{PCB} with flat readout strips. Similarly to the case of the free wires in the \ac{HD} design, these perforations allow for the drifting electrons to pass by two induction planes, before being collected. Making the comparison of simulation results easier between these two modules, the coordinate system of the \ac{VD} geometry has been rotated by \SI{90}{\degree} along the $z$--axis. This results in the $x$--axis alignment with the drift direction. Given the nature of the \ac{LArTPC} technology, in both geometries, the $y$ and $z$ coordinate reconstruction will follow from the \ac{TPC} readout, together with a time--stamp. For beam events, the remaining $x$ coordinate reconstruction follows from a simple calculation of the beam trigger and the already mentioned \ac{TPC} time--stamp. Note that for non--beam events (such as astrophysical neutrinos), the interaction vertex along the drift volume can only be reliably reconstructed from the \ac{TPC} charge for \ac{APA}-- or \ac{CPA}--crossing tracks.

Astrophysical neutrinos accessible to DUNE are three orders of magnitude lower in energy than their beam counterparts and therefore produce tracks that are very small compared to the size of the \ac{LAr} volume. To improving the reconstruction capabilities of the experiment, both detectors incorporate a \ac{PDS} which makes use of \ac{XA} modules \cite{Machado_2018} distributed over the APA (for HD) or cathode and membrane planes (for VD). These provide an additional and instantaneous light signal for any interaction which, provided a correct match from the \ac{TPC} readout is found, can be used to reconstruct the drift coordinate of low--energy interacting particles. This enables background fiducialization capabilities along the drift direction which would otherwise be impossible for any of the modules. 
%\ac{DUNE}'s active material, argon, is also unique in its sensitivity to the $\nu_e$ astrophysical flux.  The $\nu_e$~CC interaction on $^{40}$Ar has long been known to have a large scattering cross--section suitable for astroparticle measurements~\cite{raghavan1986inverse}. The PDS should allow for improvements in the low energy physics range, specifically regarding supernova and solar neutrinos ($\sim\SI{10}{\mega\electronvolt}$) and allow for a fast vertex reconstruction and improved calorimetric reconstruction of events above $\SI{5}{\mega\electronvolt}$.

%The second phase of \ac{DUNE} will include at least one LAr TPC. The TPC technology allows sub-cm tracking of events over a large volume with dynamic range suitable to precisely reconstruct neutrinos at several GeV for long-baseline neutrino oscillation measurements and astrophysicalñ neutrinos in the 5 to 50~MeV energy range. Several ideas for \ac{DUNE} phase2 have been proposed with potential to significantly lower thresholds, increasing reach of low-energy physics.  A near detector will also be constructed to mitigate systematic uncertainties related to long-baseline neutrino oscillation measurements.

\subsection{Background Model\label{sec:BACKGROUND_MODEL}}
The measurement of low--energy neutrinos (\SI{<20}{\mega\electronvolt}), particularly those from the Solar \( ^8\mathrm{B} \) chain, requires a thorough understanding of intrinsic (emitted from the \ac{LAr} or detector components) and external (penetrating from the cavern) backgrounds. One of the most studied backgrounds for \ac{LAr} detectors is Radon (\(^{222}\text{Rn}\)), which is a noble gas produced in the decay chain of \(^{238}\text{U}\)~\cite{nndc_rn222,borexino_radon}. As a gas, it can diffuse into the active volume of \ac{LAr} detectors, where its decay daughters accumulate on surfaces or remain suspended, emitting beta, gamma and alpha radiation (see Appendix~\ref{sec:APPENDIX}). Alpha particles have a very short range in \ac{LAr}, but due to their high energy and quenching, they produce significant light signals in the detector, which can be misidentified as low--energy neutrino interaction candidates in the \ac{PDS}.

Understanding the contributions from these isotopes is crucial in achieving a correct vertex and energy reconstruction required for low--energy neutrino analyses. Apart from the previously mentioned radon chain, a comprehensive list of background sources has been implemented in \ac{DUNE}'s simulation. These backgrounds can be classified into three main categories: intrinsic, component--induced, and external backgrounds from the cavern. 

Presenting \ac{DUNE}'s full background model is out of the scope of this work. Hence, we focus on providing a brief summary while discussing the features that have the biggest impact on the low--energy neutrino reconstruction analyses.

\begin{itemize}
  \item\textbf{External Backgrounds}: External radiation originates outside the detector's cryostat. The flux is divided into two kinds of particles that manage to penetrate into the detector:

        \begin{enumerate}
          \item\textbf{Cavern Neutrons}: Produced by (\(\alpha,n\)) reactions in surrounding rock or via cosmic--ray spallation, neutrons are emitted up to \SI{\sim 10}{\mega\electronvolt} in kinetic energy. They can enter the active volume where they undergo \ac{ES} until they capture producing a gamma--cascade which can mimic the topology of a \ac{CC} neutrino interaction. Some neutrons can be captured on the cryostat and other detector components contributing to the environmental particle flux with high energy gammas.
          \item\textbf{Environmental gammas}: High--energy photons from the rock or shotcrete surrounding the detector, typically originating from natural radioactivity. The highest rate contributor to this flux in \ac{DUNE}'s caver at \ac{SURF} is the shotcrete covering the walls. Low--energy gammas (\SI{< 2.7}{\mega\electronvolt}) entering the active volume have a huge flux contributing mostly to worsen the signal reconstruction by mimicking the de--excitation gammas of the \ac{CC} interaction or making the topological distinction between \ac{ES} and \ac{CC} less reliable. The higher energy external gammas (\num{2.7} $<$ E \SI{< 14}{\mega\electronvolt}), are produced by \(\gamma\)--decay of isotopes like \(^{208}\text{Tl}\) (\SI{\sim 3.5}{\mega\electronvolt}) or \(^{214}\text{Bi}\) (\SI{\sim 3.27}{\mega\electronvolt}) and secondary nuclear reactions of $(n,\gamma)$ interactions in certain geological environments (\SI{> 6}{\mega\electronvolt}). These high--energy gamma rays are highly site--specific, depending on the local neutron flux and rock composition. Therefore, in situ measurements are fundamental to assess their impact and model accurately their contribution as fake signals to the different physics analyses.
          % \item \textbf{Cosmic rays}: Even in underground environments, muons and cosmogenically produced secondaries can lead to spallation products and radiogenic activation of detector materials.
        \end{enumerate}

  \item \textbf{Component-Induced Backgrounds}: These can arise from the materials used to construct the detector itself. Material screening and selection, as well as radon mitigation systems, are essential to reduce these sources. These include decay chains of long--lived isotopes introduced during fabrication or operation. Some of the most prominent are:

        \begin{enumerate}
          \item \textbf{Uranium (\(^{238}\)U)}: Produces alpha, beta, and gamma emissions, as well as daughters like \(^{214}\)Bi and \(^{210}\)Pb.
          \item \textbf{Thorium (\(^{232}\)Th)}: Similar to the uranium chain, with prominent gamma emitters like \(^{208}\)Tl (\SI{2.6}{\mega\electronvolt}).
          \item \textbf{Potassium (\(^{40}\)K)}: Often present in plastics and glass (like the ones used for traditional photon detectors). It emits a \SI{1.46}{\mega\electronvolt} gamma line.
          \item \textbf{Radon (\(^{222}\)Rn)}: Present in atmospheric air and therefore very difficult to prevent.
        \end{enumerate}

  \item \textbf{Intrinsic Liquid Argon Backgrounds}: When talking about the properties of \ac{DUNE}'s \ac{LAr} we are generalizing the properties of (\(^{40}\)Ar) but certain radioactive isotopes are either present naturally in atmospheric argon or may be introduced as impurities:

        \begin{enumerate}
          \item \textbf{Argon (\(^{39}\)Ar)}: A beta emitter with an endpoint energy of \SI{565}{\kilo\electronvolt} and a half--life of 269 years. It dominates the low--energy spectrum in atmospheric argon with an activity of \SI{\sim1}{\becquerel\per\kilo\gram}. Together with the external gammas, this background contributes to the smearing of signals and worsens the \ac{CC} and \ac{ES} separation.
          \item \textbf{Argon (\(^{42}\)Ar) and Potassium (\(^{42}\)K)}: Rare isotopes that decay via beta emission. \(^{42}\)K emits an electron (\(Q \sim 3.5\)\unit{\mega\electronvolt}), which is particularly problematic due to its similarity to low--energy solar neutrino events.
          \item \textbf{Krypton (\(^{85}\)Kr)}: A beta emitter with a \(Q\)-value of \SI{687}{\kilo\electronvolt}, introduced via contamination from air or during purification. Often requires cryogenic distillation to remove.
          \item \textbf{Radon (\(^{222}\)Rn)}: In addition to emanation from components, radon can dissolve in liquid argon.
        \end{enumerate}
\end{itemize}

%The relative importance of these isotopes depends on the source of argon: underground argon (UAr) has significantly reduced \(^{39}\)Ar activity compared to atmospheric argon (AAr).
In \ac{DUNE}, \ac{LAr} purification and recirculation installations are planned to ensure constant operating conditions. A detailed background model informed by in situ measurements, simulations, and radiopurity assays has been developed and is under constant revision to ensure accurate predictions of the background fluxes in the \ac{LAr} active volume. The model is implemented in \ac{LArSoft}, the software framework used for \ac{DUNE}'s simulations, and is used to generate realistic event samples for detector response studies and physics analyses.

\begin{figure}[ht]
    \centering
    \includegraphics[width=.49\linewidth]{Images/ParticleID_KE_Distribution.png}
    \includegraphics[width=.49\linewidth]{Images/ParticlePDG_KE_Distribution.png}
  \caption{Simulated particle flux from background model v.3 on top of predicted solar spectrum. (left) Divided by the background component. (right) Divided by particle type.\label{fig:BACKGROUND_MODEL}}
\end{figure}

In~\cref{fig:BACKGROUND_MODEL}, we show the latest simulated particle flux from the background model (v.3.6) on top of the predicted solar spectrum as a function of the particles' kinetic energy. The left panel shows the flux divided by background components, while the right panel shows it divided by particle type. The flux is dominated by external gammas, argon and krypton at low-- and external gammas and neutrons at high--energies. These figures highlight the importance of understanding these backgrounds by highlighting the huge difference in flux between the low--energy neutrino signal and the background components. Without a proper energy calibration, the smearing of the background components can lead to a significant loss of signal efficiency and a degradation of the energy resolution.